\section{缔约过失责任}
\begin{enumerate}
  \item 从开始接触就遵循合同相对性原理。
  \item 构成要件。
        \begin{enumerate}
          \item 先合同阶段。不以缔约为目的的\zd{磋商、洽谈、谈判、商谈}。
          \item 过失（义务的违反）。诚信原则产生先合同义务。\zd{说明、注意、保护、提请注意、协助保密}。
          \item 对方\zd{信赖利益损失}。赔 \yi 订约费用。\zd{车马费、食宿费}；\er 订约机会。另行订立合同产生的\zd{差价损失}。
          \item 有因果关系。
        \end{enumerate}
  \item 具体类型判断。2 大类，6 小类，3 3 各一对。
        \begin{enumerate}
          \item 未成立时。
                \begin{enumerate}
                  \item 恶意磋商。区别正常的商业风险（正常经营风险）。\zd{戏耍对方、贻误商机、破坏交易意图}。
                  \item 欺诈。\zd{故意隐瞒}与合同订立有关的重要事实。
                  \item 侵犯商业秘密和信息。\zd{泄露、不正当使用}。
                \end{enumerate}
          \item 成立。
                \begin{enumerate}
                  \item 可撤销情形。
                  \item 侵犯商业秘密和信息。
                  \item \zd{恶意不办理登记、报批}。报批条款 \yjt 独立生效 \yjt 违约；非报批条款 \yjt 缔约过失。
                \end{enumerate}
        \end{enumerate}
\end{enumerate}
